% Part: incompleteness
% Chapter: introduction
% Section: decidability

\documentclass[../../../include/open-logic-section]{subfiles}

\begin{document}

\olfileid{inc}{int}{dec}

\olsection{અનિર્ણેયતા અને અપૂર્ણતા}

ગડેલની અપૂર્ણતાનાં પ્રમેયોની સાબિતી માટે
વાક્યરચનાનું અંકગણિતીકરણ જરૂરી છે. પરંતુ
તેના વિના પણ સિદ્ધાંત બધા !!{decidable}
સંબંધોને !!{represents} કરે છે એવી ધારણાથી
કેટલાંક સરસ પરિણામો મેળવી શકાય છે. તેની
સાબિતી અટકવાની સમસ્યાની અનિર્ણેયતાની
સાબિતી જેવી વિકર્ણ દલીલ છે.

\begin{thm}
જો $\Gamma$ એવો સુસંગત સિદ્ધાંત હોય જે
દરેક !!{decidable} સંબંધને !!{represents}
કરે છે, તો $\Gamma$ !!{decidable} નથી.
\end{thm}

\begin{proof}
ધારો કે~$\Gamma$ !!{decidable} છે. હવે
બતાવીએ કે તે દરેક !!{decidable} સંબંધને
!!{represents} કરતો હોય તો~$\Gamma$ અસંગત થવો જ
પડે.

!!^{decidable} ગુણધર્મો (એકસ્થાનિક સંબંધો)
એક મુક્ત ચલવાળાં !!{formula}s વડે નિરૂપાય
છે. આવા બધાં !!{formula}sની સંગણનીય પરિગણના
$!A_0(x)$, $!A_1(x)$, \dots, લો. હવે
નીચેનો ગણ~$D \subseteq \Nat$ વિચારો:
\[
D = \Setabs{n}{\Gamma \Proves \lnot !A_n(\num{n})}
\]
$D$ !!{decidable} છે, કારણ કે $n \in D$
છે કે નહીં તે તપાસવા પહેલાં~$!A_n(x)$
અને તેમાંથી~$\lnot !A_n(\num{n})$ ગણી
શકાય. $!A_n(x)$માં $x$ની દરેક મુક્ત
ઘટનાની જગ્યાએ પદ~$\num{n}$ મૂકવું અને
$!A(\num{n})$ પહેલાં~$\lnot$ મૂકવું સ્પષ્ટ
રીતે યાંત્રિક કામ છે. ધારણા મુજબ~$\Gamma$
!!{decidable} હોવાથી
$\lnot !A(\num{n}) \in \Gamma$ છે કે નહીં
તપાસી શકાય. હોય તો~$n \in D$ અને
ન હોય તો~$n \notin D$. આમ~$D$
પણ !!{decidable} છે.
\sourcecorrection{OLINC-004}{સ્થિર અંગ્રેજી મૂળમાં
આ ફકરાના પછીના બે સૂત્રોમાં~$!A$ લખ્યું છે;
પરંતુ અહીં પરિગણનામાંથી પસંદ થયેલું
$!A_n$ ઉદ્દિષ્ટ છે. મૂળનાં સૂત્રો જાળવી
આ સૂચન બાજુમાં આપ્યું છે.}

$\Gamma$ બધાં !!{decidable} ગુણધર્મોને
!!{represents} કરતો હોવાથી~$D$ને પણ
!!{represents} કરે છે. $\Gamma$માં~$D$ને
!!{represents}s કરનારાં !!{formula}s પણ
$!A_0(x)$, $!A_1(x)$, \dots,ની યાદીમાં
જ છે. તેથી એવો નંબર~$d$ લો કે
$!A_d(x)$, $\Gamma$માં~$D$ને
!!{represents} કરે. જો $d \notin D$ હોય,
તો $!A_d(x)$ એ~$D$ને !!{represents}
કરતું હોવાથી $\Gamma \Proves
\lnot !A_d(\num{d})$. પરંતુ તેનો અર્થ
એ કે~$d$ એ~$D$ની વ્યાખ્યાયિત શરત
સંતોષે છે, એટલે~$d \in D$. આ
$d \notin D$થી વિરોધાભાસ છે. તેથી
પરોક્ષ સાબિતીથી~$d \in D$.

$d \in D$ હોવાથી~$D$ની વ્યાખ્યા મુજબ
$\Gamma \Proves \lnot !A_d(\num{d})$.
બીજી તરફ, $!A_d(x)$ એ~$\Gamma$માં
$D$ને !!{represents} કરતું હોવાથી
$\Gamma \Proves !A_d(\num{d})$.
આથી~$\Gamma$ અસંગત છે.
\end{proof}

\begin{explain}
પહેલાનો પ્રમેય બતાવે છે કે બધાં
!!{decidable} સંબંધોને !!{represents}
કરતો કોઈ સુસંગત સિદ્ધાંત !!{decidable}
હોઈ શકે નહીં. આપણે બતાવીશું કે~$\Th{Q}$
ખરેખર બધા !!{decidable} સંબંધોને
!!{represents}s કરે છે. તેથી~$\Th{Q}$ને
સમાવતા બધા સિદ્ધાંતો, જેમ કે~$\Th{PA}$
અને~$\Th{TA}$, પણ આવું કરે છે અને તેથી
!!{decidable} નથી. (આ બધા સિદ્ધાંતો
પ્રમાણભૂત નિદર્શમાં સત્ય હોવાથી સુસંગત છે.)

આ પરિણામથી પ્રથમ અપૂર્ણતા-પ્રમેયનું એક
નબળું સ્વરૂપ પણ મળે છે. જે સિદ્ધાંત
!!{axiomatizable} અને !!{complete} બંને
હોય તે !!{decidable} હોય છે. તેથી
!!{axiomatizable} હોય અને બધા
!!{decidable} ગુણધર્મોને !!{represents}s
કરતા સુસંગત સિદ્ધાંતો !!{complete}
હોઈ શકતા નથી.
\end{explain}

\begin{thm}
જો $\Gamma$ !!{axiomatizable} અને
!!{complete} હોય, તો તે !!{decidable} છે.
\end{thm}

\begin{proof}
દરેક અસંગત સિદ્ધાંત !!{decidable} છે,
કારણ કે તે બધાં !!{sentence}s સમાવે છે.
આથી ``શું $!A \in \Gamma$?''નો જવાબ
હંમેશાં ``હા'' હોય છે; એટલે તેને નક્કી
કરી શકાય છે.

હવે ધારો કે~$\Gamma$ સુસંગત છે અને
!!{axiomatizable} તથા !!{complete} પણ છે.
!!{axiomatizable} હોવાથી~$\Gamma$
!!{computably enumerable} છે. ખરેખર,
$\Gamma$નાં સ્વયંસિદ્ધ વાક્યોમાંથી મળતી
બધી સાચી !!{derivation}sને સંગણનીય
વિધેયથી પરિગણી શકાય છે. સાચી
!!{derivation}માંથી તે જે !!{sentence}
!!{derive}s કરે છે તે ગણી શકાય છે;
એટલે આ બંનેને જોડીને~$\Gamma$ના બધા
પ્રમેયોની પરિગણના કરતું સંગણનીય વિધેય
મળે છે. $\Gamma$ સુસંગત અને !!{complete}
હોવાથી !!{sentence} એ~$\Gamma$નો પ્રમેય છે ત્યારે અને
માત્ર ત્યારે જ્યારે તેનો નિષેધ~$\lnot !A$ પ્રમેય નથી.
આથી $!A \in \Gamma$ છે કે નહીં તે
આ રીતે નક્કી કરી શકાય: $\Gamma$ના
બધા પ્રમેયોની પરિગણના કરો. યાદીમાં~$!A$
આવે ત્યારે $\Gamma \Proves !A$ જાણીએ.
યાદીમાં~$\lnot !A$ આવે ત્યારે
$\Gamma \Proves/ !A$ જાણીએ. $\Gamma$
!!{complete} હોવાથી આ બેમાંથી એક ઘટના
અંતે બને જ છે, એટલે પ્રક્રિયા જવાબ આપે છે.
\end{proof}

\begin{cor}
\ollabel{cor:incompleteness}
જો $\Gamma$ સુસંગત, !!{axiomatizable}
અને દરેક !!{decidable} ગુણધર્મને
!!{represents} કરતો હોય, તો તે
!!{complete} નથી.
\end{cor}

\begin{proof}
$\Gamma$ !!{complete} હોત તો
પહેલાના પ્રમેય મુજબ તે !!{decidable}
હોત (કારણ કે તે !!{axiomatizable}
અને સુસંગત છે). પરંતુ~$\Gamma$
દરેક !!{decidable} ગુણધર્મને
!!{represents} કરતો હોવાથી પહેલા
પ્રમેય મુજબ !!{decidable} નથી.
\end{proof}

\begin{prob}
$\Th{TA} = \Setabs{!A}{\Sat{N}{!A}}$
!!{axiomatizable} નથી તે બતાવો.
$\Th{TA}$ બધા નિર્ણેય ગુણધર્મોનું
નિરૂપણ કરે છે તે ધારી શકો છો.
\end{prob}

દાખલા તરીકે~$\Th{Q}$ બધા
!!{decidable} ગુણધર્મોને !!{represents}
કરે છે તે સ્થાપિત કર્યા પછી, આ પરિણામથી
$\Th{Q}$ અપૂર્ણ છે એમ મળે છે. પરંતુ તેની
સાબિતી સ્વતંત્ર !!{sentence}નું ઉદાહરણ
આપતી નથી; ફક્ત આવું !!a{sentence}
અસ્તિત્વ ધરાવે છે તે બતાવે છે. ચોક્કસ
ઉદાહરણ માટે વાક્યરચનાનું અંકગણિતીકરણ
કરી ગડેલની મૂળ સાબિતીનો વિચાર અનુસરવો
પડે. અને પહેલો દાવો—$\Th{Q}$ ખરેખર
બધા !!{decidable} ગુણધર્મોને
!!{represents}s કરે છે—હજી બતાવવાનો છે.

દરેક \emph{રસપ્રદ} સિદ્ધાંત અપૂર્ણ કે
અનિર્ણેય હોય એવું નથી. ઘણા સિદ્ધાંતો
રસપ્રદ ગાણિતિક તથ્યો વ્યક્ત કરી શકે
એટલા શક્તિશાળી છે, છતાં ગડેલના પરિણામની
શરતો સંતોષતા નથી. દાખલા તરીકે,
$\Th{Pres} = \Setabs{!A \in \Lang{L_{A^+}}}{\Sat{N}{!A}}$,
એટલે ગુણાકાર-ચિહ્ન~$\times$ વિનાની
અંકગણિતની ભાષાના !!{sentence}sનો
પ્રમાણભૂત નિદર્શમાં સત્ય હોય એવો ગણ,
પૂર્ણ અને નિર્ણેય બંને છે. આ સિદ્ધાંત
\emph{પ્રેસબર્ગર અંકગણિત} કહેવાય છે અને
માત્ર~$\Obj{0}$, $\prime$ અને~$+$ વડે
વ્યક્ત થઈ શકે એવાં પ્રાકૃતિક સંખ્યાઓ
વિશેનાં બધાં સત્યો સાબિત કરે છે.

\end{document}
